% GENERATED FILE. Edit canonical lessons, then run npm run generate:books.
% canonical-source-hash: fnv1a64:e0495a32ef3896c0
% canonical-lessons: FR-C274-survenir, FR-C274-derouler, FR-C274-avoir-lieu, FR-C274-preceder, FR-C274-durer

\chapter{Ce qui est arrivé}
\label{ch:fr-ce-qui-est-arrive}
\hlchaptermodality{274}

\begin{chapteropening}
\textbf{By the end of this chapter:} \emph{I can say that something occurred, unfolded, took place, came before something else or lasted.}

Complete the last lesson of chapter 274: I can say that something occurred, unfolded, took place, came before something else or lasted.
\end{chapteropening}

\section[survenir]{survenir — to occur, to arise unexpectedly}

\label{lesson:FR-C274-survenir}



\begin{quote}
Before the new one: say the French for lately, then the French for soon.
\end{quote}

\subsection*{You'll want to know: survenir}
\textbf{survenir} — \textquotedblleft{}to occur, to arise unexpectedly\textquotedblright{}.

\begin{quote}
Un problème est survenu pendant le voyage.
\end{quote}

\begin{quote}
\emph{A problem arose during the trip.}
\end{quote}

Conjugated like \emph{venir}, with \emph{être}: \emph{il est survenu}.

\begin{cousinweb}
\emph{survenir} comes from Latin \emph{supervenīre} \textquotedblleft{}to come upon, to arrive unexpectedly\textquotedblright{}, from \emph{super} \textquotedblleft{}over\textquotedblright{} and \emph{venīre} \textquotedblleft{}to come\textquotedblright{} (Lewis \& Short).
\end{cousinweb}

\begin{grammarlens}[title={What you've built}]
One more word for telling, explaining and describing at B1.
\end{grammarlens}

\subsection*{Your turn}
\begin{itemize}
\raggedright
  \item \emph{Say it:} \emph{survenir}
  \item \emph{Say it:} \emph{survenir}, once more
  \item \emph{Say it:} a sentence of your own with \emph{survenir}
  \item \emph{Recall:} say the French for lately, then the French for soon, then say \emph{survenir} again
\end{itemize}

\begin{tcolorbox}[breakable,colback=teal!4,colframe=teal!35!black,title={Before you move on}]
What does survenir mean? (\textquotedblleft{}To occur, to arise unexpectedly\textquotedblright{}.) Say it once more.
\end{tcolorbox}
\section[dérouler]{dérouler — to unfold, to take place (as se dérouler); to unroll}

\label{lesson:FR-C274-derouler}



\begin{quote}
Before the new one: say the French for soon, then the French for to occur.
\end{quote}

\subsection*{You'll want to know: dérouler}
\textbf{dérouler} — \textquotedblleft{}to unfold, to take place (as se dérouler); to unroll\textquotedblright{}.

\begin{quote}
Le concert s'est déroulé dans le parc.
\end{quote}

\begin{quote}
\emph{The concert took place in the park.}
\end{quote}

Mostly reflexive in this sense: \emph{se dérouler}. Without \emph{se}, it means to unroll a carpet or a map.

\begin{cousinweb}
\emph{dérouler} is built on \emph{rouler} \textquotedblleft{}to roll\textquotedblright{}, which comes from Old French \emph{roeler}, built on Latin \emph{rota} \textquotedblleft{}wheel\textquotedblright{} (Lewis \& Short).
\end{cousinweb}

\begin{grammarlens}[title={What you've built}]
One more word for telling, explaining and describing at B1.
\end{grammarlens}

\subsection*{Your turn}
\begin{itemize}
\raggedright
  \item \emph{Say it:} \emph{dérouler}
  \item \emph{Say it:} \emph{dérouler}, once more
  \item \emph{Say it:} a sentence of your own with \emph{dérouler}
  \item \emph{Recall:} say the French for soon, then the French for to occur, then say \emph{dérouler} again
\end{itemize}

\begin{tcolorbox}[breakable,colback=teal!4,colframe=teal!35!black,title={Before you move on}]
What does dérouler mean? (\textquotedblleft{}To unfold, to take place (as se dérouler); to unroll\textquotedblright{}.) Say it once more.
\end{tcolorbox}
\section[avoir lieu]{avoir lieu — to take place, to be held}

\label{lesson:FR-C274-avoir-lieu}



\begin{quote}
Before the new one: say the French for to occur, then the French for to unfold.
\end{quote}

\subsection*{You'll want to know: avoir lieu}
\textbf{avoir lieu} — \textquotedblleft{}to take place, to be held\textquotedblright{}.

\begin{quote}
La réunion a eu lieu hier.
\end{quote}

\begin{quote}
\emph{The meeting took place yesterday.}
\end{quote}

For planned events: \emph{la fête aura lieu samedi}.

\begin{cousinweb}
\emph{lieu} is inherited from Latin \emph{locus} \textquotedblleft{}place\textquotedblright{}, and \emph{avoir} from Latin \emph{habēre} \textquotedblleft{}to have, to hold\textquotedblright{} (Lewis \& Short): to have a place.
\end{cousinweb}

\begin{grammarlens}[title={What you've built}]
One more word for telling, explaining and describing at B1.
\end{grammarlens}

\subsection*{Your turn}
\begin{itemize}
\raggedright
  \item \emph{Say it:} \emph{avoir lieu}
  \item \emph{Say it:} \emph{avoir lieu}, once more
  \item \emph{Say it:} a sentence of your own with \emph{avoir lieu}
  \item \emph{Recall:} say the French for to occur, then the French for to unfold, then say \emph{avoir lieu} again
\end{itemize}

\begin{tcolorbox}[breakable,colback=teal!4,colframe=teal!35!black,title={Before you move on}]
What does avoir lieu mean? (\textquotedblleft{}To take place, to be held\textquotedblright{}.) Say it once more.
\end{tcolorbox}
\section[précéder]{précéder — to precede, to come before}

\label{lesson:FR-C274-preceder}



\begin{quote}
Before the new one: say the French for to unfold, then the French for to take place.
\end{quote}

\subsection*{You'll want to know: précéder}
\textbf{précéder} — \textquotedblleft{}to precede, to come before\textquotedblright{}.

\begin{quote}
Un grand silence a précédé la réponse.
\end{quote}

\begin{quote}
\emph{A long silence came before the answer.}
\end{quote}

\emph{Le jour précédent} — the day before, in a story.

\begin{cousinweb}
\emph{précéder} is a learned borrowing from Latin \emph{praecēdere} \textquotedblleft{}to go before\textquotedblright{}, from \emph{prae} \textquotedblleft{}before\textquotedblright{} and \emph{cēdere} \textquotedblleft{}to go\textquotedblright{} (Lewis \& Short).
\end{cousinweb}

\begin{grammarlens}[title={What you've built}]
One more word for telling, explaining and describing at B1.
\end{grammarlens}

\subsection*{Your turn}
\begin{itemize}
\raggedright
  \item \emph{Say it:} \emph{précéder}
  \item \emph{Say it:} \emph{précéder}, once more
  \item \emph{Say it:} a sentence of your own with \emph{précéder}
  \item \emph{Recall:} say the French for to unfold, then the French for to take place, then say \emph{précéder} again
\end{itemize}

\begin{tcolorbox}[breakable,colback=teal!4,colframe=teal!35!black,title={Before you move on}]
What does précéder mean? (\textquotedblleft{}To precede, to come before\textquotedblright{}.) Say it once more.
\end{tcolorbox}
\section[durer]{durer — to last}

\label{lesson:FR-C274-durer}



\begin{quote}
Before the new one: say the French for to take place, then the French for to precede.
\end{quote}

\subsection*{You'll want to know: durer}
\textbf{durer} — \textquotedblleft{}to last\textquotedblright{}.

\begin{quote}
Le voyage a duré trois heures.
\end{quote}

\begin{quote}
\emph{The journey lasted three hours.}
\end{quote}

\emph{Ça n'a pas duré} — it did not last.

\begin{cousinweb}
\emph{durer} is inherited from Latin \emph{dūrāre} \textquotedblleft{}to harden, to hold out, to last\textquotedblright{} (EtymDB), built on \emph{dūrus} \textquotedblleft{}hard\textquotedblright{}.
\end{cousinweb}

\begin{grammarlens}[title={What you've built}]
One more word for telling, explaining and describing at B1.
\end{grammarlens}

\subsection*{Your turn}
\begin{itemize}
\raggedright
  \item \emph{Say it:} \emph{durer}
  \item \emph{Say it:} \emph{durer}, once more
  \item \emph{Say it:} a sentence of your own with \emph{durer}
  \item \emph{Recall:} say the French for to take place, then the French for to precede, then say \emph{durer} again
\end{itemize}

\begin{tcolorbox}[breakable,colback=teal!4,colframe=teal!35!black,title={Before you move on}]
What does durer mean? (\textquotedblleft{}To last\textquotedblright{}.) Say it once more.
\end{tcolorbox}
